%==============================================================================
% TEMPLATE PROGRESS REPORT - PROYEK PENGGANTI UAS
% Mata Kuliah Kecerdasan Buatan | Teknik Komputer - Universitas Indonesia
% Semester Genap 2025/2026
%
% Dokumen ini adalah laporan kemajuan (interim) sebelum laporan akhir.
% Cakupan: Bab I (definisi masalah) - Bab III (progres pekerjaan) + Bab IV
% (pembagian tugas & deklarasi kontribusi individual).
%
% Cara pakai: sama dengan template-laporan-akhir.tex.
%==============================================================================

\documentclass[11pt,a4paper,oneside]{report}

%------ Bahasa & Font --------------------------------------------------------
\usepackage{fontspec}
\setmainfont{Times New Roman}
\setmonofont{Consolas}[Scale=0.92]

%------ Geometri & Spasi -----------------------------------------------------
\usepackage[a4paper,margin=2.5cm]{geometry}
\usepackage{setspace}
\onehalfspacing
\setlength{\parskip}{0.5em}

%------ Paket Pendukung ------------------------------------------------------
\usepackage{graphicx}
\usepackage{amsmath,amssymb,amsfonts}
\usepackage{booktabs}
\usepackage{array}
\usepackage{tabularx}
\usepackage{longtable}
\usepackage{enumitem}
\usepackage{xcolor}
\usepackage{listings}
\usepackage{hyperref}
\usepackage{titlesec}
\usepackage{tcolorbox}
\usepackage{caption}

%------ Penyesuaian Bahasa Indonesia -----------------------------------------
\renewcommand{\chaptername}{BAB}
\renewcommand{\contentsname}{DAFTAR ISI}
\renewcommand{\listfigurename}{DAFTAR GAMBAR}
\renewcommand{\listtablename}{DAFTAR TABEL}
\renewcommand{\appendixname}{LAMPIRAN}
\renewcommand{\figurename}{Gambar}
\renewcommand{\tablename}{Tabel}

%------ Format Bab -----------------------------------------------------------
\titleformat{\chapter}[display]
  {\normalfont\Large\bfseries\centering}
  {\chaptername\ \thechapter}{12pt}{\Large}
\titlespacing*{\chapter}{0pt}{-20pt}{20pt}

%------ Hyperref -------------------------------------------------------------
\hypersetup{
  colorlinks=true,
  linkcolor=black,
  urlcolor=blue,
  citecolor=black,
  pdftitle={Progress Report - MK Kecerdasan Buatan}
}

%------ Box Petunjuk Pengisian ----------------------------------------------
\newtcolorbox{petunjuk}{
  colback=yellow!10,
  colframe=orange!60!black,
  boxrule=0.5pt,
  arc=2pt,
  left=6pt,right=6pt,top=4pt,bottom=4pt,
  fontupper=\small\itshape
}

%------ Box Status (selesai / berlangsung / belum) --------------------------
\newcommand{\statusbox}[1]{\textbf{[STATUS: #1]}}

%------ Listing untuk JSON ---------------------------------------------------
\lstdefinelanguage{json}{
  basicstyle=\ttfamily\footnotesize,
  showstringspaces=false,
  breaklines=true,
  frame=single,
  rulecolor=\color{gray!50},
  string=[s]{"}{"},
  stringstyle=\color{blue!60!black}
}

%==============================================================================
\begin{document}

%==============================================================================
% HALAMAN SAMPUL
%==============================================================================
\begin{titlepage}
\centering
\vspace*{1cm}

{\Large\bfseries PROGRESS REPORT \par}
{\large PROYEK PENGGANTI UAS \par}
\vspace{0.3cm}
{\large Mata Kuliah Kecerdasan Buatan \par}
{\large Semester Genap 2025/2026 \par}

\vspace{2cm}

% Logo opsional
% \includegraphics[width=0.25\textwidth]{logo-ui.png}

\vspace{1.5cm}

{\LARGE\bfseries [JUDUL PROYEK YANG DESKRIPTIF] \par}
\vspace{0.5cm}
{\large Subjudul: Topik / Model AI yang Digunakan \par}

\vspace{2cm}

{\large Periode Laporan: \par}
{\large [DD Bulan 2026] -- [DD Bulan 2026] \par}

\vspace{1cm}

{\large Disusun oleh: \par}
\vspace{0.5cm}
\begin{tabular}{rl}
{[Nama Anggota 1]} & NPM. [xxxxxxxxxx] \\
{[Nama Anggota 2]} & NPM. [xxxxxxxxxx] \\
{[Nama Anggota 3]} & NPM. [xxxxxxxxxx] \\
% [Nama Anggota 4] & NPM. [xxxxxxxxxx] \\
\end{tabular}

\vspace{1.5cm}

{\large Dosen Pengampu: \par}
{\large [Nama Dosen Pengampu] \par}

\vfill

{\large DEPARTEMEN TEKNIK KOMPUTER \par}
{\large FAKULTAS TEKNIK \par}
{\large UNIVERSITAS INDONESIA \par}
{\large [BULAN] 2026 \par}
\end{titlepage}

%==============================================================================
% DAFTAR ISI
%==============================================================================
\tableofcontents

%==============================================================================
% RINGKASAN EKSEKUTIF
%==============================================================================
\chapter*{RINGKASAN EKSEKUTIF}
\addcontentsline{toc}{chapter}{Ringkasan Eksekutif}

\begin{petunjuk}
Maksimum 200 kata. Berikan ringkasan: (1) topik yang dipilih + masalah yang dipecahkan, (2) progres saat ini dalam 1--2 kalimat (persentase selesai relatif terhadap rencana awal).
\end{petunjuk}

[ISI DI SINI]

\vspace{0.5em}
\noindent\textbf{Persentase progres keseluruhan:} \fbox{[XX]\%}

%==============================================================================
% BAB I - PENDAHULUAN & DEFINISI MASALAH
%==============================================================================
\chapter{PENDAHULUAN \& DEFINISI MASALAH}
\label{bab:pendahuluan}

\section{Topik yang Dipilih}
\begin{petunjuk}
Sebutkan 1 dari 5 topik yang disediakan tim pengajar. Alasan pemilihan dalam 1 paragraf.
\end{petunjuk}

[ISI DI SINI]

\section{Rumusan Masalah (Tentatif)}
\begin{petunjuk}
1--3 pertanyaan riset. Tandai \textit{tentatif} jika masih mungkin berubah pada laporan akhir.
\end{petunjuk}

\begin{enumerate}[label=RM\arabic*.,leftmargin=*]
    \item {[Rumusan masalah pertama]}
    \item {[Rumusan masalah kedua]}
\end{enumerate}

\section{Ruang Lingkup \& Batasan Sementara}
\begin{petunjuk}
Apa yang dikerjakan vs tidak dikerjakan. Wajib mencantumkan batasan yang \textbf{sudah disepakati tim} hingga laporan ini ditulis.
\end{petunjuk}

[ISI DI SINI]

%==============================================================================
% BAB II - PENDEKATAN AI YANG DIRENCANAKAN
%==============================================================================
\chapter{PENDEKATAN AI YANG DIRENCANAKAN}
\label{bab:pendekatan}

\section{Studi Literatur yang Telah Dilakukan}
\begin{petunjuk}
Sebutkan paper / sumber utama yang sudah dibaca. Boleh dalam bentuk daftar singkat (penulis, tahun, judul, satu kalimat insight). Belum perlu lengkap seperti tinjauan pustaka di laporan akhir; cukup yang sudah dikaji sejauh ini.
\end{petunjuk}

\begin{enumerate}[leftmargin=*]
    \item {[Sumber 1] -- [insight singkat]}
    \item {[Sumber 2] -- [insight singkat]}
    \item {[Sumber 3] -- [insight singkat]}
\end{enumerate}

\section{Arsitektur Model yang Dipilih (Sementara)}
\begin{petunjuk}
Sebutkan model/arsitektur yang akan digunakan + varian spesifik (mis. ResNet-50 vs ResNet-18, GRU vs LSTM, BERT-base vs DistilBERT). Sertakan diagram blok jika sudah dirancang.
\end{petunjuk}

[ISI DI SINI]

\begin{figure}[h!]
\centering
\fbox{\parbox[c][3.5cm][c]{0.7\textwidth}{\centering [Diagram arsitektur sementara]}}
\caption{Sketsa arsitektur model yang direncanakan}
\label{fig:arsitektur-sementara}
\end{figure}

\section{Justifikasi Awal Pemilihan Pendekatan}
\begin{petunjuk}
Justifikasi singkat (1--2 paragraf): mengapa memilih arsitektur ini dibanding alternatif. Tidak perlu sedalam laporan akhir.
\end{petunjuk}

[ISI DI SINI]

%==============================================================================
% BAB III - PROGRES PEKERJAAN SAAT INI
%==============================================================================
\chapter{PROGRES PEKERJAAN SAAT INI}
\label{bab:progres}

\section{Ringkasan Status per Tahap}
\begin{petunjuk}
Gunakan tabel berikut untuk menggambarkan status keseluruhan. Status: $\square$ Belum mulai, $\bullet$ Sedang berjalan, $\checkmark$ Selesai.
\end{petunjuk}

\begin{table}[h!]
\centering
\caption{Status per tahapan pekerjaan}
\label{tab:status-tahapan}
\begin{tabularx}{\textwidth}{lXl}
\toprule
\textbf{Tahap} & \textbf{Aktivitas} & \textbf{Status} \\
\midrule
Pencarian dataset      & [deskripsi]       & [$\checkmark$/$\bullet$/$\square$] \\
EDA \& pra-pemrosesan  & [deskripsi]       & [$\checkmark$/$\bullet$/$\square$] \\
Implementasi model     & [deskripsi]       & [$\checkmark$/$\bullet$/$\square$] \\
Training awal          & [deskripsi]       & [$\checkmark$/$\bullet$/$\square$] \\
Evaluasi \& metrik     & [deskripsi]       & [$\checkmark$/$\bullet$/$\square$] \\
Hyperparameter tuning  & [deskripsi]       & [$\checkmark$/$\bullet$/$\square$] \\
Error analysis         & [deskripsi]       & [$\checkmark$/$\bullet$/$\square$] \\
Penulisan laporan akhir & [deskripsi]      & [$\checkmark$/$\bullet$/$\square$] \\
\bottomrule
\end{tabularx}
\end{table}

\section{Dataset \& Pra-pemrosesan}
\begin{petunjuk}
Status dataset: sudah didapat / belum / dalam proses akuisisi. Jika sudah ada: sumber, ukuran, distribusi kelas. Pra-pemrosesan yang sudah berjalan.
\end{petunjuk}

\statusbox{[selesai / berlangsung / belum]}

[ISI DI SINI]

\section{Implementasi Model}
\begin{petunjuk}
Sejauh mana kode model sudah ditulis. Modul apa yang sudah selesai, mana yang masih dalam pengembangan. Sertakan tautan komit/branch repository jika relevan.
\end{petunjuk}

\statusbox{[selesai / berlangsung / belum]}

[ISI DI SINI]

\section{Hasil Eksperimen Sementara (Jika Ada)}
\begin{petunjuk}
Jika sudah ada hasil training awal (meskipun belum tuning penuh): tampilkan metrik sementara. Tidak masalah jika belum optimal --- justru penting untuk menggambarkan trajektori. Tandai sebagai \textit{preliminary}.
\end{petunjuk}

\begin{table}[h!]
\centering
\caption{Hasil eksperimen sementara (\textit{preliminary})}
\label{tab:hasil-sementara}
\begin{tabular}{lccc}
\toprule
\textbf{Konfigurasi} & \textbf{Train} & \textbf{Val} & \textbf{Catatan} \\
\midrule
{[Run 1]} & [nilai] & [nilai] & [catatan] \\
{[Run 2]} & [nilai] & [nilai] & [catatan] \\
\bottomrule
\end{tabular}
\end{table}

[Komentar singkat terhadap hasil sementara.]

%==============================================================================
% BAB IV - PEMBAGIAN TUGAS & DEKLARASI KONTRIBUSI INDIVIDUAL
%==============================================================================
\chapter{PEMBAGIAN TUGAS \& DEKLARASI KONTRIBUSI INDIVIDUAL}
\label{bab:kontribusi}

\begin{petunjuk}
\textbf{BAB INI WAJIB DITULIS SECARA MANDIRI OLEH MASING-MASING ANGGOTA} (bukan oleh ketua kelompok). Isi Sub-bab ini akan digunakan untuk men-generate pertanyaan ujian lisan yang personal. Klaim penguasaan konsep per anggota dideklarasikan dalam berkas \texttt{metadata.json} (field \texttt{konsep\_dikuasai\_*}) yang merupakan bagian tidak terpisahkan dari laporan ini.

Gandakan template di bawah untuk setiap anggota (3--4 orang).
\end{petunjuk}

\section{[Nama Anggota 1] -- NPM. [xxxxxxxxxx]}

\subsection*{IV.1.1 Komponen yang Dikerjakan}
Daftar konkret bagian proyek yang menjadi tanggung jawab penuh saya:
\begin{itemize}[leftmargin=*]
    \item \textbf{Kode/modul:} [mis. \texttt{src/preprocessing/}, fungsi \texttt{train\_step()} di \texttt{src/train.py}]
    \item \textbf{Bagian arsitektur:} [mis. \textit{attention layer}, \textit{replay buffer}, konstruksi graf]
    \item \textbf{Bagian analisis:} [mis. \textit{confusion matrix}, \textit{ablation study}]
    \item \textbf{Bagian laporan:} [mis. Bab II.3, Bab III.4]
\end{itemize}

\subsection*{IV.1.2 Pernyataan Penggunaan AI \textit{Assistant}}
\begin{itemize}[leftmargin=*]
    \item \textbf{Tool yang dipakai:} [mis. ChatGPT 5.1, Claude Opus 4.7, GitHub Copilot --- sebutkan tool \& versi]
    \item \textbf{Untuk bagian apa:} [mis. \textit{debugging} CUDA OOM, \textit{drafting} paragraf \textit{related work}, eksplorasi ide]
    \item \textbf{Tingkat ketergantungan:} [rendah / sedang / tinggi] (estimasi $\sim$[X]\% output berasal dari AI)
    \item \textbf{Pernyataan:} Saya memahami dan dapat mempertanggungjawabkan setiap baris kode/teks yang saya klaim sebagai bagian saya.
\end{itemize}

%-------------------- Anggota 2 (gandakan blok di atas) ----------------------
\section{[Nama Anggota 2] -- NPM. [xxxxxxxxxx]}

\subsection*{IV.2.1 Komponen yang Dikerjakan}
\begin{itemize}[leftmargin=*]
    \item \textbf{Kode/modul:} [\ldots]
    \item \textbf{Bagian arsitektur:} [\ldots]
    \item \textbf{Bagian analisis:} [\ldots]
    \item \textbf{Bagian laporan:} [\ldots]
\end{itemize}

\subsection*{IV.2.2 Pernyataan Penggunaan AI \textit{Assistant}}
\begin{itemize}[leftmargin=*]
    \item \textbf{Tool yang dipakai:} [\ldots]
    \item \textbf{Untuk bagian apa:} [\ldots]
    \item \textbf{Tingkat ketergantungan:} [\ldots]
    \item \textbf{Pernyataan:} Saya memahami dan dapat mempertanggungjawabkan setiap baris kode/teks yang saya klaim sebagai bagian saya.
\end{itemize}

%-------------------- Anggota 3 ---------------------------------------------
\section{[Nama Anggota 3] -- NPM. [xxxxxxxxxx]}

\subsection*{IV.3.1 Komponen yang Dikerjakan}
\begin{itemize}[leftmargin=*]
    \item \textbf{Kode/modul:} [\ldots]
    \item \textbf{Bagian arsitektur:} [\ldots]
    \item \textbf{Bagian analisis:} [\ldots]
    \item \textbf{Bagian laporan:} [\ldots]
\end{itemize}

\subsection*{IV.3.2 Pernyataan Penggunaan AI \textit{Assistant}}
\begin{itemize}[leftmargin=*]
    \item \textbf{Tool yang dipakai:} [\ldots]
    \item \textbf{Untuk bagian apa:} [\ldots]
    \item \textbf{Tingkat ketergantungan:} [\ldots]
    \item \textbf{Pernyataan:} Saya memahami dan dapat mempertanggungjawabkan setiap baris kode/teks yang saya klaim sebagai bagian saya.
\end{itemize}

%-------------------- Anggota 4 (jika ada, hapus jika hanya 3 anggota) ------
% \section{[Nama Anggota 4] -- NPM. [xxxxxxxxxx]}
% [\ldots gandakan blok di atas \ldots]

%==============================================================================
% LAMPIRAN
%==============================================================================
\appendix
\renewcommand{\thechapter}{\Alph{chapter}}
\titleformat{\chapter}[display]
  {\normalfont\Large\bfseries\centering}
  {\appendixname\ \thechapter}{12pt}{\Large}

%------ Lampiran A: Metadata JSON (versi sementara) -------------------------
\chapter{Berkas Metadata Proyek (Versi Sementara)}
\label{lamp:metadata}

\begin{petunjuk}
Tempelkan isi berkas \texttt{metadata.json} versi saat ini. Field yang belum pasti boleh ditandai dengan \texttt{"TBD"}. Akan diperbarui di laporan akhir.
\end{petunjuk}

\begin{lstlisting}[language=json]
{
  "kode_kelompok": "K00",
  "topik_pilihan": 0,
  "judul_proyek": "...",
  "status_progress": "TBD",
  "anggota": [
    {
      "nama": "...",
      "npm": "...",
      "focus_area": ["..."]
    }
  ]
}
\end{lstlisting}

%------ Lampiran B: Tautan Repositori ---------------------------------------
\chapter{Tautan Repositori \& Sumber}
\label{lamp:tautan}

\begin{description}[leftmargin=4cm,labelwidth=3.5cm]
    \item[GitHub repo:] \url{https://github.com/...}
    \item[Google Colab:] \url{https://colab.research.google.com/...}
    \item[Sumber dataset:] \url{...}
\end{description}

\end{document}
